\section{Related Work}
\subsection{Synaptic eligibility and learning in recurrent spiking circuits}
Reward-modulated STDP and related three-factor rules connect synaptic activity to reinforcement received later. Florian derived reward-modulated spike-based rules. Izhikevich modeled dopamine acting on activity-dependent synaptic traces, and Legenstein and colleagues analyzed learning with reward-modulated STDP \cite{florian2007reinforcement,izhikevich2007distal,legenstein2008learning}. Yagishita and colleagues found a limited interval in which dopamine promotes structural plasticity after glutamatergic activation of striatal spines \cite{yagishita2014critical}. This supports a delayed interaction between activity and modulation, without identifying a particular computational eligibility variable with its biological substrate.

Reward centering can affect what a reward-modulated rule learns. Fr\'emaux and colleagues showed that inappropriate centering can leave a reward-independent contribution to the expected update for a class of Hebbian rules \cite{fremaux2010functional}. This result is specific to the analyzed rules and provides no guarantee for arbitrary reward-weighted updates. Their later spiking actor--critic model used continuous-time temporal-difference learning for control, including simulated navigation \cite{fremaux2013continuous}. Brzosko and colleagues combined slice experiments on sequential neuromodulation with a computational navigation model \cite{brzosko2017sequential}, relating the timing of plasticity modulation to behavioral learning.

E-prop trains recurrent spiking networks by pairing local synaptic factors with approximated neuron-specific feedback \cite{bellec2020}. SuperSpike also combines spike-history factors with error signals using surrogate derivatives \cite{zenke2018superspike}. Our recurrent network contains regular and adaptive LIF neurons and uses symmetric e-prop for prediction learning. The local recurrence retains adaptive-threshold derivatives and stops derivatives through recurrent spike transmission and reset. Episodic REINFORCE updates a separate motor readout \cite{williams1992}, using a surrogate derivative through the spiking motor circuit. There is no learned value function, direct reward update of recurrent synapses, or separate pair-based STDP rule.

\subsection{Astrocytic signaling, inhibition, and memory}
Neuronal modulation by astrocytes varies with the pathway and preparation. Bowser and Khakh found that ATP can depolarize hippocampal interneurons and increase inhibition onto their targets \cite{bowser2004atp}. Torres and colleagues described a route from extracellular calcium depletion through astrocytic ATP release to increased inhibitory transmission \cite{torres2012extracellular}. These studies motivate inhibitory feedback without specifying the threshold rule used here. Other pathways affect different parts of the circuit. Astrocytic activation can enhance hippocampal plasticity and memory \cite{adamsky2018astrocytic}; an endocannabinoid--astrocyte--D-serine loop affects dendritic integration and spatial learning \cite{bohmbach2022astrocytic}. Modeling this loop has linked it to plasticity changes during reversal learning \cite{squadrani2024reversal}. These mechanisms act at different sites, so evidence for one cannot be used interchangeably with evidence for another.

The diversity of calcium signaling also limits a general slow-memory interpretation. Sun and colleagues found developmental regulation of astrocytic metabotropic glutamate signaling \cite{sun2013glutamate}. De Ceglia and colleagues identified a specialized astrocyte population capable of glutamatergic gliotransmission \cite{deceglia2023specialized}. Petravicz and colleagues observed preserved performance in several behavioral assays after conditional disruption of IP$_3$R2-dependent astrocytic signaling \cite{petravicz2014astrocyte}. Calcium fluctuations nevertheless persist in astrocytic processes in IP$_3$R2-deficient mice \cite{srinivasan2015calcium}, so the manipulation does not remove all astrocytic calcium signaling. Together with rapid cortical responses \cite{stobart2018cortical}, these findings argue for assumptions tied to the modeled pathway.

Neuron--astrocyte models have implemented several forms of memory. Gordleeva and colleagues modeled the maintenance and retrieval of neuronal patterns through astrocyte-dependent modulation, followed by situation-based memory in neuron--astrocyte networks \cite{gordleeva2021working,gordleeva2025situation}. De Pitt\`a and Brunel modeled persistent, dynamic, and activity-silent working memory arising from synapse--astrocyte interactions \cite{depitta2022multiple}. These results show possible uses of astrocytic state. They leave open whether such a state is necessary for navigation when position is observed directly. Facilitation-based memory \cite{mongillo2008synaptic} and other neuronal or synaptic mechanisms remain alternative explanations for the benefit of persistence.

\subsection{Neuron--astrocyte models of adaptive behavior}
Gong and colleagues used a neuron--astrocyte model for contextual guidance in stationary and nonstationary bandit tasks \cite{gong2024astrocytes}. This work connects slow glial variables to adaptive behavior. Its tasks and modulation rules differ from our use of predicted spatial features to recruit action-specific inhibition. A computational benefit in the present controller would not identify the cellular mechanism of reward-guided behavior.

Astrocyte-regulated reinforcement learning has also been used for spiking motor control. Han and Sengupta combined astrocytic regulation with reward-modulated STDP in a central pattern generator for simulated quadruped locomotion \cite{han2026astrocyte}. In Dong and He's astrocyte-gated multi-timescale plasticity, slow astrocytic variables modulate plasticity directly \cite{dong2026agmp}. Our familiarity memory changes the excitability of inhibitory motor neurons and thereby alters action probabilities. It can also affect later prediction learning by changing the transitions collected by the policy.

The reduced glutamate--calcium cascade and calcium-dependent interneuron-threshold modulation come from Tsybina and colleagues' navigation preprint \cite{tsybina2026}. Our model queries distributed spatial features with learned predictions of candidate next states. It does not assign a glial unit to each environmental location. This separates the stored history from the recurrent representation used to predict outcomes. The comparison tests the complete prediction--memory--motor pathway; the calcium suppression motif itself is inherited.

\subsection{Auxiliary prediction and the interpretation of controls}
UNREAL showed that auxiliary tasks can improve representations shared with an RL controller \cite{jaderberg2017reinforcement}. Here, the predictor learns next positions and spatial features, and action selection uses its feature predictions. The actor reads detached recurrent activity, so reward has no direct gradient path into the recurrent network. Prediction learning and reward learning remain coupled through the transitions collected by the policy. E-prop has also been used in a spiking wavefront navigation planner \cite{espino2024rapid}, whose planning representation differs from the candidate-familiarity mechanism studied here.

The controls test distinct parts of the prediction--memory--motor pathway. Replacing calcium with an exponential trace tests whether a simpler persistent state can drive the same inhibition. Independent parameter selection compares complete systems, including their memory dynamics, gain, and optimization settings. Removing memory at evaluation tests its contribution with learned weights fixed. Shared-input probes assess transition prediction without relying on the trajectories selected by each policy. These controls leave unanswered what the same actor would have learned without memory throughout training.
